\loesung{

The code for this programming exercise can be found in the \texttt{hw\_shapley\_py\_sol.ipynb} or \texttt{hw\_shapley\_SHAP\_py\_sol.ipynb} files for Python and in the \texttt{hw\_shapley\_R\_sol.Rmd} or \texttt{hw\_shapley\_SHAP\_R\_sol.Rmd} files for R.

\begin{enumerate}[a)]
	\item Pseudocode for predicting the Man of the Match probability through a random forest
	
	\begin{algorithm}[H]
		\caption{Man of the Match}
		\begin{algorithmic}[1]
			\State \texttt{df} $\gets$ read in \textit{fifa.csv}
			\State \texttt{df}['Man of the Match'] $\gets$ replace 'Yes' by TRUE (else FALSE)
			\State \texttt{df} $\gets$ adapt \texttt{df} if needed for random forest model (e.g. removing NAs, removing non-integer features)
			\State \texttt{train, test} $\gets$ split \texttt{df} in train and test data
			\State \texttt{classifier\_RF} $\gets$ random forest fit using \texttt{train}
		\end{algorithmic}
	\end{algorithm}

    For specific calculation results see the implemented examples, because they depend on the specific implementation, the seed and other parameters. 

	\item Pseudocode of \texttt{marginal\_vfunc(S, observation, X, predict, nr\_samples)}
	
	\begin{algorithm}[H]
		\caption{\texttt{marginal\_vfunc()}}
		\begin{algorithmic}[1]
			\Require \texttt{S}: feature index or indices, or feature names
			\Require \texttt{observation}: observation where to evaluate the marginal expectation
			\Require \texttt{X}: feature matrix
			\Require \texttt{predict}: function returning the predictions of the RF model
			\Require \texttt{nr\_samples}: number of samples
			\State \texttt{X\_tmp} $\gets$ sample \texttt{nr\_samples} samples from \texttt{X} (with replacing)
			\State \texttt{X\_tmp} $\gets$ replace features \texttt{S} with respective values from \texttt{observation}
			\State \textbf{return} mean of \texttt{predict(X\_tmp)}
		\end{algorithmic}
	\end{algorithm}

    \item For specific results, see the implementations. We plug \texttt{marginal\_vfunc()} in as the value function when calling \texttt{shapley\_perm\_approx()} and need to take care that the other arguments of \texttt{marginal\_vfunc()} are handled correctly.

	\item Pseudocode of \texttt{shap\_weights()}
	
	\begin{algorithm}[H]
		\caption{\texttt{shap\_weights()}}
		\begin{algorithmic}[1]
			\Require \texttt{mask}: (binary) coalition feature space
			\State \texttt{p} $\gets$ number of features in \texttt{mask}
			\State \texttt{zs} $\gets$ coalition sizes (row sums of \texttt{mask})
			\State \texttt{numerator} $\gets$ \texttt{p} - 1
			\State \texttt{denominator} $\gets$ (binomial coefficient of \texttt{p} over \texttt{zs}) $*$ \texttt{zs} $*$ (\texttt{p} - \texttt{zs})
			\State \textbf{return} \texttt{numerator} / \texttt{denominator}
		\end{algorithmic}
	\end{algorithm}

    \newpage

	Pseudocode of \texttt{replace\_dataset()}
	
	\begin{algorithm}[H]
		\caption{\texttt{replace\_dataset()}}
		\begin{algorithmic}[1]
			\Require \texttt{obs}: observation
			\Require \texttt{X}: feature matrix
			\Require \texttt{nr\_samples}: number of samples
			\State \texttt{X\_new} $\gets$ sample \texttt{nr\_samples} samples from \texttt{X} (with replacing)
			\State \texttt{obs\_rep} $\gets$ matrix with \texttt{nr\_samples} rows containing \texttt{obs} in each row
			\State \texttt{mask} $\gets$ \texttt{nr\_samples} $\times$ \texttt{p} matrix with entries drawn from $\text{Bernoulli}(0.5)$
			\State \texttt{X\_new} $\gets$ replace entries where mask equals 1 with entry from \texttt{obs}
			\State \textbf{return} \texttt{X\_new}, \texttt{mask}
		\end{algorithmic}
	\end{algorithm}

	Pseudocode of \texttt{shap\_data()}
	
	\begin{algorithm}[H]
		\caption{\texttt{shap\_data()}}
		\begin{algorithmic}[1]
			\Require \texttt{obs}: observation
			\Require \texttt{X}: feature matrix
			\Require \texttt{nr\_samples}: number of samples
			\Require \texttt{predict}: prediction model
			\State \texttt{X\_new, mask} $\gets$ \texttt{replace\_dataset(obs, X, nr\_samples)}
			\State \texttt{weight} $\gets$ \texttt{shap\_weights(mask)}
			\State \texttt{pred} $\gets$ \texttt{predict(X\_new)}
			\State \textbf{return} \texttt{mask, pred, weight}
		\end{algorithmic}
	\end{algorithm}

	Pseudocode of \texttt{kernel\_shap()}
	
	\begin{algorithm}[H]
		\caption{\texttt{kernel\_shap()}}
		\begin{algorithmic}[1]
			\Require \texttt{obs}: observation
			\Require \texttt{X}: feature matrix
			\Require \texttt{nr\_samples}: number of samples
			\Require \texttt{predict}: prediction model
			\State \texttt{mask, pred, weight} $\gets$ \texttt{shap\_data(obs, X, nr\_samples, predict)}
			\State \texttt{lm} $\gets$ weighted linear regression model using \texttt{mask, pred} and \texttt{weight}
			\State \textbf{return} coefficients and intercept of \texttt{lm}
		\end{algorithmic}
	\end{algorithm}

    \textit{Note:} This simplified implementation uses unconstrained weighted least squares. The original KernelSHAP algorithm enforces the efficiency constraint $\phi_0 = \E[\fh(\Xv)]$ and $\sum_j \phi_j = \fh(\xv) - \phi_0$ via infinite weights on the empty and full coalitions. Our sampling-based version may not exactly satisfy this property.

	\item
    Here an example, see also in the implementations. Again, specific results and accuracy depend on the specific implementation.

    \texttt{kernel\_shap(X\_test[1, ], X\_train, 1000, classifier\_RF)}

    The resulting coefficients are the SHAP values for the respective features. A positive SHAP value means that the feature pushes the prediction for this observation towards the positive class (Man of the Match = Yes), while a negative value pushes it towards the negative class. The intercept represents the base value (average prediction). The magnitudes indicate the relative importance of each feature for this particular prediction.
	\begin{table}[H]
		%\scriptsize
		\centering
		\begin{tabular}{rrrr}
			\hline 
			(Intercept) &Goal.Scored & Ball.Possession..  & Attempts \\

			0.52535421 & 0.06126913 & -0.04780005 & -0.09366948  \\
			\hline
			On.Target& Off.Target & Blocked &  Corners  \\

			-0.01463136 & -0.04102268 & -0.02460521 & -0.04113441  \\
			\hline
			Offsides & Free.Kicks & Saves & Pass.Accuracy..  \\

			0.04542043 & -0.03651103 & -0.05663537 & -0.01282451   \\
			\hline
			Passes & Distance.Covered..Kms. & Fouls.Committed & Yellow.Card   \\

			0.01454292 & -0.02847718 & -0.02017669 & 0.08108944  \\
			\hline
			Yellow...Red & Red & Goals.in.PSO &   \\

			-0.00132795 & -0.05331953 & 0.00916541 & \\
			\hline
		\end{tabular} 
	\end{table}	
\end{enumerate} 
}
